% Source: https://artofproblemsolving.com/community/c3h2297476_parabola
% Stable identity: aops-topic-2297476-q01
\textbf{(a)} Factor $a$ out of the terms with $x$ and complete the square:
\[y=a\left(x^2+\dfrac bax\right)+c=a\left(x+\dfrac{b}{2a}\right)^2-\dfrac{b^2}{4a}+c.\]
So $y=a(x-h)^2+k$ with
\[h=-\dfrac{b}{2a},\qquad k=c-\dfrac{b^2}{4a}=\dfrac{4ac-b^2}{4a}.\]
Since $(x-h)^2\ge0$, with equality only at $x=h$: if $a>0$, then $y\ge k$ for every point of the graph, with equality only at $(h,k)$, and $y$ grows without bound as $x$ moves away from $h$; so $(h,k)$ is the lowest point and the parabola opens upward. If $a<0$, then $y\le k$, $(h,k)$ is the highest point and the parabola opens downward. Replacing $x-h$ by $-(x-h)$ does not change $y$, so the graph is symmetric about the vertical line $x=h$ through the vertex.

\textbf{(b)} The equation $x=ay^2+by+c$ is the equation of (a) with the letters $x$ and $y$ exchanged, so its graph is the reflection of the graph in (a) in the line $y=x$ (a parabola with a horizontal axis). Completing the square in $y$,
\[x=a\left(y+\dfrac{b}{2a}\right)^2+c-\dfrac{b^2}{4a},\]
so the vertex is $\left(c-\dfrac{b^2}{4a},\,-\dfrac{b}{2a}\right)$. If $a>0$, every point has $x\ge c-\dfrac{b^2}{4a}$: the curve lies to the right of the vertex and opens to the right; if $a<0$ it opens to the left. (So the two ``vertex forms'' are $y=a(x-h)^2+k$, opening up or down, and $x=a(y-k)^2+h$, opening right or left, both with vertex $(h,k)$.)

\textbf{(c)} $y=2x^2-12x+13=2(x^2-6x)+13=2(x-3)^2-18+13=2(x-3)^2-5$: the vertex is $\boxed{(3,-5)}$ and, since $a=2>0$, the parabola \textbf{opens upward}.

$x=-y^2+4y+1=-(y^2-4y)+1=-(y-2)^2+4+1=-(y-2)^2+5$: the vertex is $\boxed{(5,2)}$ and, since $a=-1<0$, the parabola \textbf{opens to the left}.

\begin{asy}
usepackage("amsmath");
size(12cm);
defaultpen(fontsize(10pt));
import graph;
real f(real x){return 2*(x-3)^2-5;}
draw((-0.8,0)--(6,0),Arrow(6)); draw((0,-6)--(0,3),Arrow(6));
label("$x$",(6,0),E); label("$y$",(0,3),N);
draw(graph(f,1.1,4.9),blue+linewidth(1));
draw((3,-6)--(3,2.5),dashed+gray);
dot((3,-5),blue); label("$(3,-5)$",(3,-5),E,blue);
label("$y=2x^2-12x+13$",(3,-6.2),S,blue+fontsize(9pt));
picture q;
pair g(real t){return (-(t-2)^2+5,t);}
draw(q,(-2.4,0)--(6.6,0),Arrow(6)); draw(q,(0,-1.2)--(0,5.6),Arrow(6));
label(q,"$x$",(6.6,0),E); label(q,"$y$",(0,5.6),N);
draw(q,graph(g,-0.55,4.55),red+linewidth(1));
draw(q,(-2.4,2)--(6.2,2),dashed+gray);
dot(q,(5,2),red); label(q,"$(5,2)$",(5,2),E,red);
label(q,"$x=-y^2+4y+1$",(2,-1.4),S,red+fontsize(9pt));
add(shift(12,-2)*q);
\end{asy}
